\section{Conclusion}
Ce chapitre a défini le cadre méthodologique de l'étude comparative entre l'agriculture de précision et l'agriculture traditionnelle. Le dispositif repose sur deux zones de trois plots de 2~m\textsuperscript{2} (gombo, navet, laitue) : une zone test, pilotée à l'aide du capteur 8-en-1 et du DHT22, et une zone témoin, irriguée selon un calendrier fixe. Les besoins en eau sont estimés par la méthode FAO-56 Penman-Monteith ($ET_c = ET_0 \times K_c$) et les besoins en nutriments par le calcul du déficit en N, P et K. Enfin, la zone test a bénéficié d'un amendement du sol (biochar, poudre de coque d'arachide, compost, engrais granulé bio) pour lutter contre la salinité. Ce protocole constitue la base de la mise en œuvre sur le terrain, présentée au chapitre suivant.
